% ===========================================================================
%  NIELS BOHR — Atomteori og naturbeskrivelse
%  Filosofiske Skrifter, Bind I. Kobenhavn: Forlaget Philosophia, 2013.\ Genoptryk af Kjobenhavns Universitets festskrift 1929 (J.H. Schultz, 1958).
%
%  Danish text of Volume I of an eight-volume edition
%  (four Danish, four English) of Bohr's philosophical writings.
%
%  WHAT IS NEW IN THIS EDITION
%  Both paginations are marked in the margins, so a passage can be cited
%  either way:
%      [n]  left margin, bold   — page n of the ORIGINAL publication
%                                 (journal, proceedings, Festskrift)
%      |n|  right margin, italic — page n of the COLLECTED volume
%                                 (Philosophia 2013)
%  Every essay's own original printing is named in its \provenance line and
%  in the header comment of its fragment under parts/da/.
%
%  STRUCTURE
%  This file is a shell. The text lives in one .texfrag per essay under
%  parts/da/, \input in reading order below. Fragments carry
%  .texfrag, not .tex, so the root Makefile does not try to build them
%  standalone (they have no preamble).
%
%  SOURCES ON DISK (~/bibliotek/Bohr, Niels)
%    1.1-atomteori-mekanik-1925.pdf
%    1.2-quantum-postulate-1927.pdf
%    1.3-virkningskvantet-1929.pdf
%    1.4-atomteorien-grundprincipperne-1929.pdf
%
%  EDITORIAL RECORD — keep this header in step with note.md.
%  Verified page map:        (none yet)
%  Errata applied:           (none — diplomatic transcription; errors stand)
%  Normalisation:            (to be settled at first batch)
%  Emphasis conventions:     (to be settled at first batch)
%
%  STATUS: skeleton. 5 essays scaffolded, 0 transcribed.
% ===========================================================================

\documentclass[12pt,a4paper]{article}

\def\bohrlang{danish}
\def\bohrvol{I}
\def\bohrtitle{Atomteori og naturbeskrivelse}
\def\bohrsubtitle{Filosofiske Skrifter, Bind I. K\o{}benhavn: Forlaget Philosophia, 2013.\\ Genoptryk af Kj\o{}benhavns Universitets festskrift 1929 (J.H.\ Schultz, 1958).}
\def\bohrrunning{Atomteori og naturbeskrivelse}

\input{../bohr-preamble.texfrag}

\title{\textbf{Atomteori og naturbeskrivelse}\\[8pt]
  \large Filosofiske Skrifter · Bind I}
\author{Niels Bohr}
\date{Filosofiske Skrifter, Bind I. K\o{}benhavn: Forlaget Philosophia, 2013.\\ Genoptryk af Kj\o{}benhavns Universitets festskrift 1929 (J.H.\ Schultz, 1958).}

\begin{document}
\maketitle
\thispagestyle{empty}

\clearpage
\tableofcontents

% ---------------------------------------------------------------------------
%  STANDALONE OFFPRINT -- generated by bohr-edition/build.py.
%  Edit 01-indledende-oversigt.texfrag, not this file.
%  The essay also appears in ../../transcription.pdf, the complete volume.
% ---------------------------------------------------------------------------
\documentclass[12pt,a4paper]{article}

\def\bohrlang{danish}
\def\bohrvol{I}
\def\bohrtitle{Indledende Oversigt}
\def\bohrsubtitle{Atomteori og naturbeskrivelse}
\def\bohrrunning{Atomteori og naturbeskrivelse}
\newif\ifoffprint
\offprinttrue

\input{../../../bohr-preamble.texfrag}

\begin{document}
\thispagestyle{fancy}
{\footnotesize\itshape Særtryk af: Niels Bohr, \emph{Atomteori og naturbeskrivelse}, Bind~I
 (Forlaget Philosophia, K\o{}benhavn 2013).\par}
\medskip

% ---------------------------------------------------------------------------
%  STANDALONE OFFPRINT -- generated by bohr-edition/build.py.
%  Edit 01-indledende-oversigt.texfrag, not this file.
%  The essay also appears in ../../transcription.pdf, the complete volume.
% ---------------------------------------------------------------------------
\documentclass[12pt,a4paper]{article}

\def\bohrlang{danish}
\def\bohrvol{I}
\def\bohrtitle{Indledende Oversigt}
\def\bohrsubtitle{Atomteori og naturbeskrivelse}
\def\bohrrunning{Atomteori og naturbeskrivelse}
\newif\ifoffprint
\offprinttrue

\input{../../../bohr-preamble.texfrag}

\begin{document}
\thispagestyle{fancy}
{\footnotesize\itshape Særtryk af: Niels Bohr, \emph{Atomteori og naturbeskrivelse}, Bind~I
 (Forlaget Philosophia, K\o{}benhavn 2013).\par}
\medskip

% ---------------------------------------------------------------------------
%  STANDALONE OFFPRINT -- generated by bohr-edition/build.py.
%  Edit 01-indledende-oversigt.texfrag, not this file.
%  The essay also appears in ../../transcription.pdf, the complete volume.
% ---------------------------------------------------------------------------
\documentclass[12pt,a4paper]{article}

\def\bohrlang{danish}
\def\bohrvol{I}
\def\bohrtitle{Indledende Oversigt}
\def\bohrsubtitle{Atomteori og naturbeskrivelse}
\def\bohrrunning{Atomteori og naturbeskrivelse}
\newif\ifoffprint
\offprinttrue

\input{../../../bohr-preamble.texfrag}

\begin{document}
\thispagestyle{fancy}
{\footnotesize\itshape Særtryk af: Niels Bohr, \emph{Atomteori og naturbeskrivelse}, Bind~I
 (Forlaget Philosophia, K\o{}benhavn 2013).\par}
\medskip

\input{01-indledende-oversigt.texfrag}
\end{document}

\end{document}

\end{document}

\input{parts/da/02-atomteori-og-mekanik.texfrag}
% ---------------------------------------------------------------------------
%  STANDALONE OFFPRINT -- generated by bohr-edition/build.py.
%  Edit 03-kvantepostulatet.texfrag, not this file.
%  The essay also appears in ../../transcription-da.pdf, the complete volume.
% ---------------------------------------------------------------------------
\documentclass[12pt,a4paper]{article}

\def\bohrlang{danish}
\def\bohrvol{I}
\def\bohrtitle{Kvantepostulatet og Atomteoriens seneste Udvikling}
\def\bohrsubtitle{Atomteori og naturbeskrivelse}
\def\bohrrunning{Atomteori og naturbeskrivelse}
\newif\ifoffprint
\offprinttrue

\input{../../../bohr-preamble.texfrag}

\begin{document}
\thispagestyle{fancy}
{\footnotesize\itshape Særtryk af: Niels Bohr, \emph{Atomteori og naturbeskrivelse}, Bind~I
 (Forlaget Philosophia, K\o{}benhavn 2013).\par}
\medskip

% ---------------------------------------------------------------------------
%  STANDALONE OFFPRINT -- generated by bohr-edition/build.py.
%  Edit 03-kvantepostulatet.texfrag, not this file.
%  The essay also appears in ../../transcription-da.pdf, the complete volume.
% ---------------------------------------------------------------------------
\documentclass[12pt,a4paper]{article}

\def\bohrlang{danish}
\def\bohrvol{I}
\def\bohrtitle{Kvantepostulatet og Atomteoriens seneste Udvikling}
\def\bohrsubtitle{Atomteori og naturbeskrivelse}
\def\bohrrunning{Atomteori og naturbeskrivelse}
\newif\ifoffprint
\offprinttrue

\input{../../../bohr-preamble.texfrag}

\begin{document}
\thispagestyle{fancy}
{\footnotesize\itshape Særtryk af: Niels Bohr, \emph{Atomteori og naturbeskrivelse}, Bind~I
 (Forlaget Philosophia, K\o{}benhavn 2013).\par}
\medskip

% ---------------------------------------------------------------------------
%  STANDALONE OFFPRINT -- generated by bohr-edition/build.py.
%  Edit 03-kvantepostulatet.texfrag, not this file.
%  The essay also appears in ../../transcription-da.pdf, the complete volume.
% ---------------------------------------------------------------------------
\documentclass[12pt,a4paper]{article}

\def\bohrlang{danish}
\def\bohrvol{I}
\def\bohrtitle{Kvantepostulatet og Atomteoriens seneste Udvikling}
\def\bohrsubtitle{Atomteori og naturbeskrivelse}
\def\bohrrunning{Atomteori og naturbeskrivelse}
\newif\ifoffprint
\offprinttrue

\input{../../../bohr-preamble.texfrag}

\begin{document}
\thispagestyle{fancy}
{\footnotesize\itshape Særtryk af: Niels Bohr, \emph{Atomteori og naturbeskrivelse}, Bind~I
 (Forlaget Philosophia, K\o{}benhavn 2013).\par}
\medskip

\input{03-kvantepostulatet.texfrag}
\end{document}

\end{document}

\end{document}

\input{parts/da/04-virkningskvantet.texfrag}
% ---------------------------------------------------------------------------
%  STANDALONE OFFPRINT -- generated by bohr-edition/build.py.
%  Edit 05-grundprincipperne.texfrag, not this file.
%  The essay also appears in ../../transcription.pdf, the complete volume.
% ---------------------------------------------------------------------------
\documentclass[12pt,a4paper]{article}

\def\bohrlang{danish}
\def\bohrvol{I}
\def\bohrtitle{Atomteorien og Grundprincipperne for Naturbeskrivelsen}
\def\bohrsubtitle{Atomteori og naturbeskrivelse}
\def\bohrrunning{Atomteori og naturbeskrivelse}
\newif\ifoffprint
\offprinttrue

\input{../../../bohr-preamble.texfrag}

\begin{document}
\thispagestyle{fancy}
{\footnotesize\itshape Særtryk af: Niels Bohr, \emph{Atomteori og naturbeskrivelse}, Bind~I
 (Forlaget Philosophia, K\o{}benhavn 2013).\par}
\medskip

% ---------------------------------------------------------------------------
%  STANDALONE OFFPRINT -- generated by bohr-edition/build.py.
%  Edit 05-grundprincipperne.texfrag, not this file.
%  The essay also appears in ../../transcription.pdf, the complete volume.
% ---------------------------------------------------------------------------
\documentclass[12pt,a4paper]{article}

\def\bohrlang{danish}
\def\bohrvol{I}
\def\bohrtitle{Atomteorien og Grundprincipperne for Naturbeskrivelsen}
\def\bohrsubtitle{Atomteori og naturbeskrivelse}
\def\bohrrunning{Atomteori og naturbeskrivelse}
\newif\ifoffprint
\offprinttrue

\input{../../../bohr-preamble.texfrag}

\begin{document}
\thispagestyle{fancy}
{\footnotesize\itshape Særtryk af: Niels Bohr, \emph{Atomteori og naturbeskrivelse}, Bind~I
 (Forlaget Philosophia, K\o{}benhavn 2013).\par}
\medskip

% ---------------------------------------------------------------------------
%  STANDALONE OFFPRINT -- generated by bohr-edition/build.py.
%  Edit 05-grundprincipperne.texfrag, not this file.
%  The essay also appears in ../../transcription.pdf, the complete volume.
% ---------------------------------------------------------------------------
\documentclass[12pt,a4paper]{article}

\def\bohrlang{danish}
\def\bohrvol{I}
\def\bohrtitle{Atomteorien og Grundprincipperne for Naturbeskrivelsen}
\def\bohrsubtitle{Atomteori og naturbeskrivelse}
\def\bohrrunning{Atomteori og naturbeskrivelse}
\newif\ifoffprint
\offprinttrue

\input{../../../bohr-preamble.texfrag}

\begin{document}
\thispagestyle{fancy}
{\footnotesize\itshape Særtryk af: Niels Bohr, \emph{Atomteori og naturbeskrivelse}, Bind~I
 (Forlaget Philosophia, K\o{}benhavn 2013).\par}
\medskip

\input{05-grundprincipperne.texfrag}
\end{document}

\end{document}

\end{document}


\end{document}
